\section{Experiment}
\vspace{-12pt}
% \subsection{Datasets and Implementation Details.}
\begin{figure}[ht]
  \centering
  % \vspace{-10pt}
  \includegraphics[width=1\linewidth]{images/experiment.pdf}
  \caption{\textbf{Multi-resolution rendering of 3D scenes.} 
    We visualize rendered images and depth maps at three resolutions: $\times1$ ({\color{blue!70}blue} box), $\times2$ ({\color{green!60!black}green} box), and $\times4$ ({\color{red!70}red} box). 
    As resolution increases, artifacts in the underlying 3D representation become more evident. 
    In the image space, they appear as dark regions caused by unfilled gaps, where hollow areas are rendered as black pixels.
    % these artifacts manifest as darkened regions due to unfilled gaps—i.e., hollow areas rendered as black pixels. 
    In the depth space, they appear as unnatural yellow regions, indicating incorrect depth predictions caused by geometric discontinuities or sparsity. Note that yellow corresponds to near surfaces and blue denotes distant regions in depth map visualization.
    % artifacts are reflected as unnatural yellow regions, which indicate incorrect depth predictions—often caused by geometric discontinuities or sparsity. 
    % Note: in the depth map, yellow corresponds to near surfaces, while blue represents distant regions.
    }
  \label{fig:experiment}
  \vspace{-15pt}
\end{figure}


\noindent\textbf{Datasets.}
% \paragraph{Datasets.} 
To evaluate our method, we follow the experimental setup in PixelSplat~\cite{charatan2024pixelsplat} and conduct experiments on the RealEstate10K (RE10K)~\cite{zhou2018stereo} and ACID~\cite{liu2021infinite} datasets. 
RE10K mainly consists of indoor real estate videos, whereas ACID contains outdoor scenes captured by aerial drones. Both datasets provide precomputed camera poses and we adhere to the official train-test splits used in prior work. 
Additionally, we evaluate our method on the DTU~\cite{jensen2014large} dataset following MVSplat~\cite{chen2024mvsplat}, on DL3DV~\cite{ling2024dl3dv} following DepthSplat~\cite{xu2024depthsplat}, and further extend our evaluation to the challenging ScanNet~\cite{dai2017scannet} dataset.

\noindent\textbf{Evaluation Metrics.}
% \paragraph{Evaluation Metrics.}
We evaluate novel view synthesis quality using standard metrics: PSNR, SSIM, and LPIPS. 
To better evaluate geometric fidelity, we additionally report high-resolution rendering consistency (HRRC) results at $2\times$ and $4\times$ resolution.

\noindent\textbf{Baselines.} 
% \paragraph{Baselines.} 
We compare our method to state-of-the-art sparse-view generalizable methods for novel view synthesis, including PixelSplat~\cite{charatan2024pixelsplat}, MVSplat~\cite{chen2024mvsplat}, TranSplat~\cite{zhang2025transplat}, HiSplat~\cite{tang2024hisplat}, and DepthSplat~\cite{xu2024depthsplat}. 
Among these, PixelSplat and HiSplat generate multiple Gaussians per pixel, while MVSplat, TranSplat, and DepthSplat predict a single Gaussian per pixel. 
Since using more primitives generally improves performance, we focus our core comparisons on the latter group to ensure a fair comparison.

\noindent\textbf{Implementation Details.} 
% \paragraph{Implementation Details.} 
Our method is implemented using PyTorch~\cite{paszke2019pytorch} and optimized using AdamW~\cite{loshchilov2017decoupled} with a cosine learning rate schedule. 
We conduct experiments with different monocular backbones from Depth Anything V2~\cite{yang2024depth} (ViT-S, ViT-B, ViT-L), referred to as Ours-S, Ours-B, and Ours-L respectively. 
We train our models for a total of 4800K iterations on an NVIDIA A100 GPU following DepthSplat~\cite{xu2024depthsplat},. 
% For the small model (Ours-S), we train for 300K iterations with a batch size of 16; for the base and large models (Ours-B/L), we train for 600K iterations with a batch size of 8. 
For the small model (Ours-S), we train for 300K iterations with a batch size of 16, while the base and large models (Ours-B and Ours-L) are trained for 600K iterations with a batch size of 8.
We adopt the encoder settings from DepthSplat~\cite{xu2024depthsplat}, but use a lower learning rate of $2 \times 10^{-6}$ for the pretrained Depth Anything V2 backbone. All other layers are trained with a learning rate of $2 \times 10^{-4}$. 
The opacity threshold $\tau_{\text{opa}}$ is set to 0.6, and the alpha normalization threshold $\tau_{\alpha}$ is set to 0.1 during training and 0.001 during evaluation. Predicted scale multipliers are clamped to the range $[\frac{1}{3}, 3]$. We train our models at $256 \times 256$ resolution for fair comparison unless otherwise specified. Furthermore, we explore higher-resoluton training at $256 \times 448$ and demonstrate the results in the appendix~\ref{hq}.

\subsection{Main Results}
\vspace{-5pt}

\begin{table}[ht]
\centering
\caption{Novel view synthesis performance on the RealEstate10k dataset.}
\label{re10k}
\fontsize{7}{11.5}\selectfont
\setlength{\tabcolsep}{4pt}
\begin{tabular}{lccc|ccc|ccc|ccc}
\toprule
& \multicolumn{3}{c|}{256$\times$256 (\textbf{Standard})} & \multicolumn{3}{c|}{512$\times$512 (\textbf{HRRC})} & \multicolumn{3}{c|}{1024$\times$1024 (\textbf{HRRC})} & \multicolumn{3}{c}{Average} \\
Method & PSNR↑ & SSIM↑ & LPIPS↓ & PSNR↑ & SSIM↑ & LPIPS↓ & PSNR↑ & SSIM↑ & LPIPS↓ & PSNR↑ & SSIM↑ & LPIPS↓ \\
\midrule
pixelSplat  & 26.049 & 0.862 & 0.137 & 25.782 & 0.868 & 0.207 & 24.920 & 0.877 & 0.269 & 25.584 & 0.869 & 0.204 \\
HiSplat  & 27.193 & 0.882 & 0.117 & 25.269 & 0.870 & 0.198 & 24.262 & 0.878 & 0.248 & 25.575 & 0.877 & 0.188 \\

\midrule
MVSplat & 26.359 & 0.868 & 0.129 & 20.408 & 0.809 & 0.290 & 17.966 & 0.755 & 0.425 & 21.578 & 0.811 & 0.281 \\
TranSplat & 26.687 & 0.875 & 0.125 & 20.610 & 0.815 & 0.286 & 18.154 & 0.761 & 0.427 & 21.817 & 0.817 & 0.279 \\
DepthSplat & \underline{27.504} & \underline{0.890} & \textbf{0.112} & 20.031 & 0.774 & 0.341 & 16.385 & 0.635 & 0.491 & 21.307 & 0.766 & 0.315 \\
\textbf{Ours-S} & 27.001 & 0.883 & 0.118 & 25.989 & \underline{0.860} & 0.223 & 24.535 & 0.835 & 0.325 & 25.842 & 0.859 & 0.222 \\
\textbf{Ours-B} & 27.447 & \underline{0.890} & \underline{0.113} & \underline{26.280} & \textbf{0.866} & \underline{0.218} & \underline{24.744} & \underline{0.838} & \underline{0.322} & \underline{26.157} & \underline{0.865} & \underline{0.217} \\
\textbf{Ours-L} & \textbf{27.537} & \textbf{0.892} & \textbf{0.112} & \textbf{26.331} & \textbf{0.866} & \textbf{0.217} & \textbf{24.897} & \textbf{0.842} & \textbf{0.320} & \textbf{26.255} & \textbf{0.867} & \textbf{0.216} \\
\bottomrule
\end{tabular}
\vspace{-15pt}
\end{table}


\begin{table}[ht]
\centering
\caption{Novel view synthesis performance on the ACID dataset.}
\label{acid}
\fontsize{7}{11.5}\selectfont
\setlength{\tabcolsep}{4pt}
\begin{tabular}{lccc|ccc|ccc|ccc}
\toprule
& \multicolumn{3}{c|}{256$\times$256 (\textbf{Standard})} & \multicolumn{3}{c|}{512$\times$512 (\textbf{HRRC})} & \multicolumn{3}{c|}{1024$\times$1024 (\textbf{HRRC})} & \multicolumn{3}{c}{Average} \\
Method & PSNR↑ & SSIM↑ & LPIPS↓ & PSNR↑ & SSIM↑ & LPIPS↓ & PSNR↑ & SSIM↑ & LPIPS↓ & PSNR↑ & SSIM↑ & LPIPS↓ \\
\midrule
pixelSplat  & 28.284 & 0.842 & 0.146 & 27.687 & 0.848 & 0.243 & 26.462 & 0.858 & 0.343 & 27.478 & 0.849 & 0.244 \\
HiSplat  & 28.737 & 0.853 & 0.132 & 25.376 & 0.833 & 0.246 & 23.988 & 0.841 & 0.314 & 25.700 & 0.842 & 0.231 \\
\midrule
MVSplat & 28.202 & \underline{0.842} & 0.145 & 17.802 & 0.711 & 0.406 & 14.784 & 0.572 & 0.567 & 20.263 & 0.708 & \underline{0.373} \\
TranSplat  & \textbf{28.337} & \textbf{0.845} & \textbf{0.143} & \underline{17.911} & \underline{0.716} & \underline{0.402} & \underline{14.956} & \underline{0.582} & \underline{0.558} & \underline{20.401} & \underline{0.714} & \underline{0.373} \\
\textbf{Ours} & \underline{28.336} & \textbf{0.845} & \underline{0.144} & \textbf{26.868} & \textbf{0.814} & \textbf{0.281} & \textbf{21.253} & \textbf{0.690} & \textbf{0.457} & \textbf{25.486} & \textbf{0.783} & \textbf{0.294} \\
\bottomrule
\end{tabular}
\vspace{-12pt}
\end{table}

\begin{table}[t]
\centering
\caption{Cross datasets performance.}
\vspace{-10pt}
\label{cross}
\fontsize{7}{11.5}\selectfont
\setlength{\tabcolsep}{4pt}
\begin{tabular}{lccc|ccc|ccc|ccc}
\toprule
& \multicolumn{3}{c|}{Scannet} & \multicolumn{3}{c|}{DL3DV} & \multicolumn{3}{c|}{DTU} & \multicolumn{3}{c}{Average} \\
Method & PSNR↑ & SSIM↑ & LPIPS↓ & PSNR↑ & SSIM↑ & LPIPS↓ & PSNR↑ & SSIM↑ & LPIPS↓ & PSNR↑ & SSIM↑ & LPIPS↓ \\
\midrule
pixelSplat  & 19.606 & 0.714 & 0.324 & 27.201 & 0.882 & 0.104 & 12.752 & 0.329 & 0.639 & 19.853 & 0.642 & 0.356 \\
HiSplat  & 19.095 & 0.691 & 0.342 & 26.242 & 0.869 & 0.112 & 16.019 & 0.671 & 0.277 & 20.452 & 0.744 & 0.244 \\
\midrule
MVSplat & 18.725 & 0.692 & 0.333 & 23.841 & 0.768 & 0.156 & 13.914 & 0.470 & 0.386 & 18.827 & 0.643 & 0.292 \\
TranSplat  & 18.944 & 0.705 & 0.332 & 23.913 & 0.771 & 0.161 & \underline{14.956} & \textbf{0.527} & \textbf{0.327} & 19.271 & 0.668 & \underline{0.273} \\
DepthSplat  & \underline{20.201} & \textbf{0.735} & \textbf{0.305} & \textbf{28.141} & \textbf{0.905} & \textbf{0.083} & 14.592 & 0.425 & 0.436 & \underline{20.978} & \underline{0.688} & 0.275 \\
\textbf{Ours} & \textbf{20.305} & \underline{0.731} & \underline{0.313} & \underline{27.384} & \underline{0.890} & \underline{0.106} & \textbf{15.544} & \underline{0.488} & \underline{0.329} & \textbf{21.078} & \textbf{0.703} & \textbf{0.249} \\
\bottomrule
\end{tabular}
\vspace{-10pt}
\end{table}

\begin{table}[ht]
\centering
\caption{Ablations study on various components.}
% \textbf{FAB}: Forced Alpha Blending; \textbf{SCP}: Surface Continuity Prior.}
\vspace{-10pt}
\label{ablation}
\fontsize{7}{11.5}\selectfont
\setlength{\tabcolsep}{4pt}
\begin{tabular}{lccc|ccc|ccc|ccc}
\toprule
& \multicolumn{3}{c|}{256$\times$256 (\textbf{Standard})} & \multicolumn{3}{c|}{512$\times$512 (\textbf{HRRC})} & \multicolumn{3}{c|}{1024$\times$1024 (\textbf{HRRC})} & \multicolumn{3}{c}{Average} \\
Method & PSNR↑ & SSIM↑ & LPIPS↓ & PSNR↑ & SSIM↑ & LPIPS↓ & PSNR↑ & SSIM↑ & LPIPS↓ & PSNR↑ & SSIM↑ & LPIPS↓ \\
\midrule
w/o FAB, SCP & 26.925 & 0.880 & 0.120 & 21.549 & 0.805 & 0.307 & 18.563 & 0.716 & 0.422 & 22.346 & 0.800 & 0.283 \\
w/o FAB & 26.481 & 0.873 & 0.128 & 21.042 & 0.776 & 0.345 & 17.576 & 0.662 & 0.474 & 21.700 & 0.770 & 0.316 \\
Full  & \textbf{27.001} & \textbf{0.883} & \textbf{0.118} & \textbf{25.989} & \textbf{0.860} & \textbf{0.223} & \textbf{24.535} & \textbf{0.835} & \textbf{0.325} & \textbf{25.842} & \textbf{0.859} & \textbf{0.222} \\
\bottomrule
\end{tabular}
\vspace{-18pt}
\end{table}

\noindent\textbf{Reconstruction Quality.}
% \paragraph{Reconstruction Quality.}
% \paragraph{3D Scene Quality.}
We report quantitative comparison on the RE10K dataset in Table~\ref{re10k} and on the ACID dataset in Table~\ref{acid}. 
Our proposed SurfSplat method consistently outperforms previous state-of-the-art methods across various metrics and datasets, especially under high-resolution rendering settings. As shown in Figure~\ref{fig:experiment}, we visualize the predicted 3D scenes rendered into both RGB and depth maps at the original, $\times2$, and $\times4$ resolutions. 
While previous methods appear visually plausible at the original resolution, their reconstructions manifest spatial inconsistencies at higher resolutions, including holes and surface gaps.
These artifacts reveal the limitations of previous feedforward 3DGS models in capturing sub-pixel-level geometry. 
Notably, DepthSplat, despite using the same encoder backbone as our method, fails to generate coherent geometry or consistent surface details, which highlights the effectiveness of our surface continuity prior and forced alpha blending strategy.

\noindent\textbf{Cross-Dataset Generalization.}
% \paragraph{Cross-Dataset Generalization.}
To assess cross-dataset generalization, we train our model on RE10K and directly conduct evaluation on DTU, DL3DV, and ScanNet datasets. 
As shown in Table~\ref{cross}, SurfSplat maintains strong performance and generalizes better than previous methods across all target domains. This demonstrates the robustness of our learned geometric prior and the general applicability of our representation even under domain shift.


\subsection{Ablation and Analysis}
% TODO: without both FAB and SCP (w/o FAB,SCP)
\begin{wrapfigure}{l}{0.5\linewidth}
    \vspace{-15pt}
    \includegraphics[width=\linewidth]{images/ablation_new.pdf}
    % \vspace{-5pt}
    \captionsetup{belowskip=-16pt}
    \caption{\textbf{Ablation study: Visualization of reconstructed 3D scenes.} 
    % From top to bottom: full model, without forced alpha blending (w/o FAB), without surface continuity prior (w/o SCP), and without both FAB and SCP (w/o FAB,SCP). 
    Our full model yields continuous and coherent surfaces, while ablated variants exhibit visible artifacts and spatial inconsistencies.}
    \label{fig:ablation}
\end{wrapfigure}

We conduct extensive ablation studies to further validate the effectiveness of key components. Specifically, we evaluate variants without both of \textit{forced alpha blending} and \textit{surface continuity prior} (denoted as \textbf{w/o FAB,SCP}) , and without \textit{forced alpha blending} (denoted as \textbf{w/o FAB}). Quantitative results are reported in Table~\ref{ablation}, and we also rendered the reconstructed 3D scenes onto three orthogonal planes in Figure~\ref{fig:ablation} to provide qualitative comparisons. Our full model yields continuous and coherent surfaces, while ablated variants exhibit visible artifacts and spatial inconsistencies.

\noindent\textbf{Surface Continuity Prior.} 
% \paragraph{Surface Continuity Prior.}  
To evaluate the impact of the surface continuity prior, we train a variant that independently predicts all Gaussian attributes without geometric coupling. Interestingly, this variant still achieves competitive novel view synthesis (NVS) performance at the original resolution, despite producing visually noisy and discontinuous surfaces. This observation highlights a key limitation of conventional NVS metrics and underscores the value of our proposed HRRC metric, which drops significantly when surface continuity is not enforced.

\noindent\textbf{Forced Alpha Blending.} 
% \paragraph{Forced Alpha Blending.} 
We also train a variant with the surface continuity prior but without forced alpha blending. We observe a clear spatial misalignment across views, as the model tends to produce fully opaque Gaussians, which occlude background information and hinder correct 3D alignment. This leads to a substantial drop in both standard and HRRC metrics.


\subsection{Effectiveness of HRRC metric}
To empirically validate the effectiveness of HRCC on native high-resolution data, we conducted additional experiments on the high-resolution version of the DL3DV dataset. We randomly sampled a representative subset for evaluation and ensured that all methods were tested under identical conditions. The results are reported in Table~\ref{table:native_hq}. Across these experiments, the relative performance rankings remained fully consistent with those observed under HRRC evaluation, even without any bicubic upsampling. This indicates that the conclusions drawn from HRRC reliably transfer to native high-resolution evaluations.

\subsection{Normal and Mesh Comparison}
Since our method naturally predicts a surface orientation for each 2DGS, we additionally generate the corresponding normal maps and reconstructed meshes to further demonstrate the effectiveness of SurfSplat. We provide a comparison with DepthSplat~\cite{yang2024depth} in Figure~\ref{fig:normal_mesh}. From this comparison, we observe that our method produces more geometrically consistent results, highlighting the improved geometric coherence induced by the surface continuity prior.


\begin{figure}
    \centering
    \includegraphics[width=\linewidth]{images/normal_mesh.pdf}
    \caption{Normal and mesh comparison with DepthSplat.}
    \label{fig:normal_mesh}
\end{figure}


\begin{table}
    \centering
    \caption{Quantitative performance comparison on high-resolution DL3DV dataset.}
    \label{table:native_hq}
    % \fontsize{7}{11.5}\selectfont
    % \setlength{\tabcolsep}{4pt}
   \begin{tabular}{l|cc|cccc}
    \toprule
    Metric & pixelSplat & HiSplat & MVSplat & TransSplat & DepthSplat & Ours \\
    \midrule
    PSNR $\uparrow$   & \underline{24.082} & 22.780 & 17.966 & 19.545 & 16.066 & \textbf{24.411} \\
    SSIM $\uparrow$   & 0.755  & \underline{0.765}  & 0.645  & 0.679  & 0.600  & \textbf{0.788} \\
    LPIPS $\downarrow$& 0.250  & \underline{0.237}  & 0.301  & 0.257  & 0.424  & \textbf{0.252} \\
    \bottomrule
\end{tabular}

\end{table}
